\section{Prevalence Arguments}
\label{appendix:prevalence}

This appendix proves Proposition~\ref{prop:e-submersion-finite-attractors}. The
proof has three steps. First, we split roots of $\QTC$ into continuous interior
and boundary root problems. Second, we use parametric transversality to show
that these root sets are discrete for prevalent dynamics. Third, we combine
discreteness with compact containment from Assumption~\ref{ass:kkq} to obtain
finiteness.

\begin{remark}
    With $C^2$ boundary, the normal vector $\vn:\partial X \to \R^d$ is well-defined and $C^1$. So we can write $\QTC$ as:
    \[
        \QTC(\vx)=
        \begin{cases}
            \vQ(\vx) & \vx \in \inter X \\
            \QTCP(\vx) & \vx \in \partial X
        \end{cases}
    \]
    where, for $\vx\in\partial X$,
    \[
        \QTCP(\vx)=
        \vQ(\vx) - \langle \vQ(\vx), \vn(\vx) \rangle_+ \vn(\vx),
    \]
    and the $+$ subscript denotes the positive part (maximum of the argument and zero).
\end{remark}

The final step of the argument uses the following compactness fact.

\begin{lemma}[Closed discrete sets in compact sets are finite]\label{lemma:closed-discrete-finite}
Let $Y$ be a topological space, and let $S\subseteq Y$ be closed in $Y$ and
discrete in $Y$. If there is a compact set $K\subseteq Y$ such that
$S\subseteq K$, then $S$ is finite.
\end{lemma}
\begin{proof}
Since $S$ is closed in $Y$, $S\cap K=S$ is closed in $K$. Hence $S$ is compact
in the subspace topology. Since $S$ is discrete, for each $s\in S$ there is an
open set $U_s\subseteq Y$ such that $U_s\cap S=\{s\}$. The collection
$\{U_s\cap S:s\in S\}$ is an open cover of the compact space $S$ by singletons,
so it has a finite subcover. Hence $S$ is finite.
\end{proof}

In Section~\ref{sec:model}, we showed that $g_{\vQ}(X)$ is finite iff $\QTC$
has finitely many roots. The useful fact about roots of a continuous function is
that they form a closed set. Although $\QTC$ is not continuous, its roots are
contained in the union of the roots of two continuous maps.

\begin{lemma}\label{lemma:break-into-two}
	   \[\mR(\QTC) \subseteq \mR(\vQ) \cup \mR(\QTCP)\]
	   where $\mR(\vQ)$ is closed in $X$ and $\mR(\QTCP)$ is closed in $\partial X$.
\end{lemma}
\begin{proof}
Take any $\vx\in \mR(\QTC)$, i.e. $\QTC(\vx)=0$. If $\vx\in \inter X$, then $\QTC(\vx)=\vQ(\vx)$, so $\vQ(\vx)=0$ and hence $\vx\in \mR(\vQ)$. If instead $\vx\in \partial X$, then $\vx\in \mR(\QTCP)$ by definition. Therefore $\mR(\QTC) \subseteq \mR(\vQ) \cup \mR(\QTCP)$.

As $\vQ$ is continuous, $\mR(\vQ)$ is closed. Under Assumption~\ref{ass:kkq}, $X$ has $C^2$ boundary. For $\vx\in\partial X$,
\[\QTCP(\vx) = \vQ(\vx) - \langle \vQ(\vx), \vn(\vx) \rangle_+ \vn(\vx)\]
As the boundary is $C^2$, $\vn$ is $C^1$, making $\QTCP$ continuous on
$\partial X$ and $\mR(\QTCP)$ closed in $\partial X$.
\end{proof}

What remains is having a condition that implies discreteness of these root sets.
Once we have discreteness, Assumption~\ref{ass:kkq} provides compact containment
and Lemma~\ref{lemma:closed-discrete-finite} gives finiteness. When a function
$f$ is $C^1$, a sufficient condition for roots of $f$ to be isolated is
transversality of $f$ to $0$: $f \pitchfork 0$. As transversality works well
with prevalence, we phrase discreteness in terms of transversality to $0$. This
works directly for $\vQ$, which is $C^1$. The boundary map $\QTCP$, however, is
not $C^1$ because it involves $\langle \vQ(\vx), \vn(\vx) \rangle_+$.

To avoid the nonsmooth positive-part term in $\QTCP$, we encode boundary roots
using the following $C^1$ map:
\[
\begin{aligned}
    \tilde{\vQ}:\partial X \times [0,\infty) \to \R^{d} \\
    \tilde{\vQ}(\vx, \alpha) = \vQ(\vx) - \alpha \vn(\vx)
\end{aligned}
 \]
The following lemma shows why $\tilde{\vQ} \pitchfork 0$ implies discreteness of $\mR(\QTCP)$.

\begin{lemma}\label{lemma:discrete-preserves}
    If $\tilde{\vQ} \pitchfork 0$, then $\mR(\QTCP)$ is discrete in $\partial X$.
\end{lemma}
\begin{proof}
    If $\tilde{\vQ} \pitchfork 0$, then $\mR(\tilde{\vQ})$ is discrete in
    $\partial X\times[0,\infty)$. Moreover, $\vx\in\mR(\QTCP)$ iff
    \[
    (\vx,\alpha(\vx))\in\mR(\tilde{\vQ}),
    \qquad
    \alpha(\vx):=\langle \vQ(\vx),\vn(\vx)\rangle_+.
    \]

    Suppose for contradiction that $\mR(\QTCP)$ is not discrete in $\partial X$.
    Then there exists a sequence of distinct points $\vx_n\in \mR(\QTCP)$ and a limit point
    $\vx^*=\lim_{n\to\infty} \vx_n\in \mR(\QTCP)$.
    By continuity of $\alpha(\cdot)$,
    $(\vx_n,\alpha(\vx_n))\to(\vx^*,\alpha(\vx^*))$. Each pair lies in
    $\mR(\tilde{\vQ})$, and the pairs are distinct because the $\vx_n$ are
    distinct. Thus $\mR(\tilde{\vQ})$ has an accumulation point, contradicting
    discreteness.
    Therefore $\mR(\QTCP)$ is discrete.
\end{proof}

So to prove discreteness of $\mR(\vQ)$ and $\mR(\QTCP)$, it suffices to prove $\vQ \pitchfork 0$ and $\tilde{\vQ} \pitchfork 0$. The following lemma connects prevalence and transversality.

\begin{restatable}[Parametric transversality from a submersion]{lemma}{lemParametricSubmersion}\label{lem:parametric_submersion}
Let $Z$ be a finite-dimensional $C^{r}$ manifold, possibly with boundary, and let
$P$ be a finite-dimensional real vector space (identified with $\R^{m}$, with
Lebesgue measure $\lambda_{P}$). Let $\vF:P \times Z\to \R^{n}$ be a $C^{r}$ map
with $r>\max(\dim Z-n,0)$ such that $\vF$ is a submersion, i.e. for every $(\vp,z)\in P\times Z$,
\[
D\vF_{(\vp,z)}:T_{\vp}P\times T_{z}Z\to \R^{n}
\quad\text{is surjective.}
\]
For each $\vp \in P$ write $\vF_{\vp}(z):=\vF(\vp,z)$.
Then for $\lambda_{P}$-a.e.\ $\vp\in P$ we have $\vF_{\vp}\pitchfork 0$.
\end{restatable}

\begin{remark}\label{remark:parametric-regularity}
    In particular, if $\dim Z=n$ then
    $r>\max(\dim Z-n,0)$ reduces to $r>0$, so $C^{1}$ regularity suffices. This
    is the typical case for the applications in this paper; when $\dim Z>n$ the
    stronger $C^{r}$ assumption above is required to apply Sard's theorem to
    the projection used in the proof (see \cite{GuilleminPollack1974}).
\end{remark}

The idea is to apply Sard's theorem to the projection from the zero set
$M:=\vF^{-1}(0)$ onto the parameter space $P$. Regular values of this projection
are exactly the parameters for which the fiber map $\vF_{\vp}$ is transverse to
$0$.

\begin{proof}\label{proof:parametric_submersion}
Let $q:=\dim Z$. Since $\vF$ is a submersion, $0$ is a regular value of $\vF$.
Using the preimage theorem, $M:=\vF^{-1}(0)\subset P\times Z$ is a
$\dim(P)+q-n$-dimensional $C^{1}$ submanifold (possibly with boundary) with
$T_M = \operatorname{ker}(D\vF)$.

Let $\pi:P\times Z\to P$ be the projection $\pi(\vp,z)=\vp$ and set
$\pi^{M}:=\pi|_{M}$. Intuitively, values of $\pi^{M}$ are $\vp$ such that for
some $z$, $\vF(\vp,z)=0$. By Sard's theorem, the set of critical values of
$\pi^{M}$ has $\lambda_{P}$-measure zero. We show that regular values of
$\pi^{M}$ correspond exactly to $\vp$ such that $\vF_{\vp}\pitchfork 0$.

\[
\begin{aligned}
\ker D \pi^{M}_{(\vp,z)}
&=\{\vvec\in T_{(\vp,z)}M:D \pi^{M}_{(\vp,z)}(\vvec)=0\}\\
&=\{(\vvec_{\vp},\vvec_{z})\in T_{(\vp,z)}M:\vvec_{\vp}=0\}\\
&=\{(0,\vvec_{z})\in T_{(\vp,z)}M\}\\
&=\{(0,\vvec_{z})\in\ker(D\vF_{(\vp,z)})\}\\
&=\{(0,\vvec_{z}):D_{z}\vF_{(\vp,z)}(\vvec_{z})=0\}\\
&=\{(0,\vvec_{z}):\vvec_{z}\in\ker(D_{z}\vF_{(\vp,z)})\}.
\end{aligned}
\]

So $\dim(\ker D \pi^{M}_{(\vp,z)}) = \dim(\ker(D_{z} \vF_{(\vp,z)}))$.
Using the rank-nullity theorem on $D \pi^{M}_{(\vp,z)}$, we have:
\[
\begin{aligned}
&\rank(D \pi^{M}_{(\vp,z)})
+\dim(\ker D \pi^{M}_{(\vp,z)})\\
&\qquad=\dim(T_{(\vp,z)}M).
\end{aligned}
\]
\[
\begin{aligned}
&\Rightarrow \rank(D \pi^{M}_{(\vp,z)})
+\dim(\ker D \pi^{M}_{(\vp,z)})\\
&\qquad=\dim(P)+q-n.
\end{aligned}
\]
Doing the same on $D_{z} \vF$, we have:
\[
\begin{aligned}
\rank(D_{z} \vF_{(\vp,z)})+\dim(\ker D_{z} \vF_{(\vp,z)})
&=\dim(Z)\\
&=q.
\end{aligned}
\]
Subtracting the two equations cancels the kernels as they are equal giving us:
\[
\rank(D \pi^{M}_{(\vp,z)})-\rank(D_{z} \vF_{(\vp,z)})
=\dim(P)-n.
\]

Now if $\vp$ is a regular value of $\pi^{M}$, then
$\rank(D \pi^{M}_{(\vp,z)}) = \dim(P)$. Therefore
$\rank(D_{z} \vF_{(\vp,z)}) = n$, which is equivalent to
$\vF_{\vp} \pitchfork 0$.

On the other hand, assume $\vF_{\vp} \pitchfork 0$. Then
$\rank(D_{z} \vF_{(\vp,z)}) = n$. Using the previous equation, we have:
\[\rank(D \pi^{M}_{(\vp,z)}) = \dim(P)\]
So $D \pi^{M}_{(\vp,z)}$ is surjective, and hence $\vp$ is a regular value of $\pi^{M}$.
\end{proof}

\begin{remark}
    When $n=\dim Z$, $\vF_{\vp} \pitchfork 0$ means $0$ is a regular value of $\vF_{\vp}$. Using the preimage theorem, we have $\dim(\vF_{\vp}^{-1}(0)) = \dim Z - n = 0$. This implies $\vF_{\vp}^{-1}(0)$, i.e. the roots of $\vF_{\vp}$, are discrete in $Z$.
\end{remark}

To obtain transversality prevalently, we perturb the dynamics along a
finite-dimensional probe space.

\begin{definition}
    For a finite-dimensional subspace $P$ of $\mathcal{Q}$, for any $\vQ \in \mathcal{Q}$, define the following perturbed evaluation maps $e^{\vQ}:P\times X\to\R^{d}$ and $e^{\tilde{\vQ}}:P \times \partial X \times [0,\infty) \to \R^{d}$:
    \[
    \begin{aligned}[t]
        e^{\vQ}(\vp,\vx)&:=(\vQ+\vp)(\vx)\\
        e^{\tilde{\vQ}}(\vp,\vx,\alpha)&:=(\vQ+\vp)(\vx) - \alpha \vn(\vx)
    \end{aligned}
    \]
\end{definition}

We can use this to prove the following sufficient condition for prevalence of transversality.

\begin{theorem}[Prevalence via a submersion probe]\label{theorem:prevalence_submersion_probe}
Let $X$ have $C^2$ boundary. Let $\mathcal Q$ be a complete metric vector space of continuously differentiable functions from $X$ to $\R^d$ and let $P\subset \mathcal Q$ be a finite-dimensional subspace.

If for all $\vQ\in\mathcal Q$, the maps $e^{\vQ}$ and $e^{\tilde{\vQ}}$ are $C^{1}$ submersions, then for prevalent $\vQ \in \mathcal{Q}$, $\mR(\vQ)$ is discrete in $X$ while $\mR(\QTCP)$ is discrete in $\partial X$.
\end{theorem}

\begin{proof}
Fix any base dynamics $\vQ_0\in\mathcal Q$. Applying
Lemma~\ref{lem:parametric_submersion} to $e^{\vQ_0}$ gives a full-measure set of
perturbations $\vp\in P$ such that $(\vQ_0+\vp)(\cdot)\pitchfork 0$, and hence
$\mR(\vQ_0+\vp)$ is discrete in $X$. Applying the same lemma to
$e^{\tilde{\vQ}_0}$ gives a full-measure set of perturbations $\vp\in P$ such
that
\[
(\vx,\alpha)\mapsto(\vQ_0+\vp)(\vx)-\alpha\vn(\vx)
\]
is transverse to $0$. By Lemma~\ref{lemma:discrete-preserves},
$\mR((\vQ_0+\vp)^{\Pi_{\TC_X}}_{|\partial X})$ is discrete in $\partial X$ for
these perturbations. Intersecting the two full-measure sets gives full measure
in $P$. Since this holds for every base $\vQ_0\in\mathcal Q$, the set of
dynamics for which both root sets are discrete is prevalent.
\end{proof}

\begin{corollary}\label{cor:prevalent-finite-roots}
Assume $X$ satisfies Assumption~\ref{ass:convex-closed}, and let $\mathcal Q$
satisfy Assumption~\ref{ass:kkq}. If there exists a finite-dimensional subspace
$P\subseteq\mathcal Q$ such that, for every $\vQ\in\mathcal Q$, the perturbed
evaluation maps $e^{\vQ}$ and $e^{\tilde{\vQ}}$ are $C^1$ submersions, then for
prevalent $\vQ\in\mathcal Q$, $\QTC$ has finitely many roots.
\end{corollary}
\begin{proof}
By Theorem~\ref{theorem:prevalence_submersion_probe}, for prevalent
$\vQ\in\mathcal Q$, $\mR(\vQ)$ is discrete in $X$ and $\mR(\QTCP)$ is discrete
in $\partial X$. By Lemma~\ref{lemma:break-into-two},
\[
\mR(\QTC)\subseteq \mR(\vQ)\cup \mR(\QTCP).
\]
The two root sets on the right are closed in their respective domains by
Lemma~\ref{lemma:break-into-two}. Assumption~\ref{ass:kkq} places them inside
compact sets. Lemma~\ref{lemma:closed-discrete-finite} therefore implies that
both $\mR(\vQ)$ and $\mR(\QTCP)$ are finite. Hence $\mR(\QTC)$ is finite.
\end{proof}

\phantomsection\label{proof:e-submersion-finite-attractors}
\propESubmersionFiniteAttractors*
\begin{proof}[Proof of Proposition~\ref{prop:e-submersion-finite-attractors}]
The hypotheses of Proposition~\ref{prop:e-submersion-finite-attractors} are
exactly the hypotheses of Corollary~\ref{cor:prevalent-finite-roots}. Hence,
for prevalent $\vQ\in\mathcal Q$, the projected dynamics $\QTC$ has finitely
many roots.

By Lemma~\ref{lemma:containment},
\[
g_{\vQ}(X)\subseteq \mR(\QTC)\cup\{\dagger\}.
\]
Thus $g_{\vQ}(X)$ is contained in a finite set, so $|g_{\vQ}(X)|<\infty$.
Therefore, for prevalent $\vQ\in\mathcal Q$, the induced dynamic solver has
finitely many reachable terminal outcomes.
\end{proof}
